\section{Realisability: an impedance ladder}
\label{sec:realizability}

UHM's substrate-closure theorem (T-153) proves the logical architecture runs
on any substrate that meets three conditions --- (C1) a finite-dimensional,
decoherence-free effective state space; (C2) CPTP (Lindblad) dynamics; (C3)
at least seven non-commuting observable modes. There is therefore no unique
machine; there is a \emph{family}, and the members differ only in how well
their physical primitives match the logical structure. We rank them by
\emph{impedance} --- the mismatch a substrate must paper over --- against four
matching criteria:

\begin{itemize}[leftmargin=1.5em,itemsep=0pt]
  \item[\textbf{(M1)}] \emph{in-place evolution} --- state advances without a
        fetch--execute cycle;
  \item[\textbf{(M2)}] \emph{Fano-wired interconnect} --- seven triple-buses,
        not a crossbar;
  \item[\textbf{(M3)}] \emph{dissipative--regenerative native dynamics} ---
        the physics is already Lindblad$+$homeostasis;
  \item[\textbf{(M4)}] \emph{seven non-commuting modes} (C3).
\end{itemize}

\begin{table}[H]
\centering\small
\caption{The realisability ladder. Higher rungs match the logical structure
with lower impedance (fewer emulation layers); lower rungs are buildable
today. The claim is \status{C}: the theorem fixes the logical architecture
and the conditions C1--C3, not the winning substrate.}
\label{tab:ladder}
\begin{tabular}{@{}p{4.3cm}cccc p{3.0cm}@{}}
\toprule
\textbf{Substrate} & \textbf{M1} & \textbf{M2} & \textbf{M3} & \textbf{M4} & \textbf{Status} \\
\midrule
Driven-dissipative analogue / photonic / open-quantum 7-mode device & \checkmark & \checkmark & \checkmark & \checkmark & lowest impedance; research \\
Memristive Fano crossbar (analogue in-memory MAC on 7 buses) & \checkmark & \checkmark & partial & \checkmark & near-term \\
Neuromorphic spiking core, Fano-topology, event-driven & \checkmark & \checkmark & partial & \checkmark & near-term \\
GPU / CPU with \DM{} in L1 and Fano-sparse kernels & emul. & emul. & emul. & \checkmark & \textbf{exists (\SYNARC{})} \\
\bottomrule
\end{tabular}
\end{table}

\begin{figure}[H]
\centering
\begin{tikzpicture}[font=\scriptsize, >=Stealth]
  % ladder rails
  \draw[very thick, talossteel] (0,0) -- (1.4,4.9);
  \draw[very thick, talossteel] (3.2,0) -- (4.6,4.9);
  % rungs (y positions along the rails)
  \foreach \t/\lab/\stat in {%
    0.06/{\textbf{Rung 0 --- GPU/CPU emulation (exists):} \DM{} in L1, Fano-sparse kernel}/{correctness, telemetry},
    0.36/{\textbf{Rung 1 --- memristive Fano crossbar:} update in the array}/{first test of B1},
    0.66/{\textbf{Rung 2 --- neuromorphic Fano core:} event-driven}/{near-term},
    0.96/{\textbf{Rung 3 --- driven-dissipative device:} Lindblad native}/{approaches the floor}}{
    \draw[thick, talosamber] ({0+1.4*\t},{4.9*\t}) -- ({3.2+1.4*\t},{4.9*\t});
  }
  \node[anchor=west] at (4.9,0.28)  {\textbf{Rung 0 --- GPU/CPU emulation (exists).} correctness + full telemetry; still moves data};
  \node[anchor=west] at (5.3,1.75)  {\textbf{Rung 1 --- memristive Fano crossbar.} update in the array; first honest test of B1};
  \node[anchor=west] at (5.7,3.22)  {\textbf{Rung 2 --- neuromorphic Fano core.} event-driven, in-place};
  \node[anchor=west] at (6.1,4.68)  {\textbf{Rung 3 --- driven-dissipative device.} physics $=$ the tick; near the energy floor};
  % arrows/axes
  \draw[->, thick, plosgray] (-0.7,0.1) -- node[left, align=center]{\tiny impedance\\ \tiny falls} (-0.1,4.7);

\end{tikzpicture}
\caption{The realisability ladder as a ladder. Every rung runs the identical
logical architecture (T-153) and the identical software; climbing removes
emulation layers (impedance falls, left arrow) and strengthens the energy
claim rung by rung. Rung numbers
match the roadmap of \S\ref{sec:discussion}.}
\label{fig:ladderfig}
\end{figure}

The bottom rung is not hypothetical: it is the current \SYNARC{} realisation,
which holds \DM{} ($784$~bytes) in L1 cache, advances it with a
Strang-split Fano-sparse kernel at $\sim\!700$--$1400$~FLOPs per tick and
arithmetic depth $\log_2 7\approx 2.8$, and reproduces every invariant and
exact law bit-for-bit against an independent oracle. It emulates M1--M3 in
software (highest impedance) and therefore does not realise the energy floor
of \S\ref{sec:energy} --- a GPU still moves data --- but it proves the logical
architecture is correct and runnable today. Each higher rung removes an
emulation layer: a memristive Fano crossbar performs the coherence update
\emph{in} the array (M1, M2 native); a driven-dissipative analogue device has
Lindblad dynamics as its \emph{native} physics (M3 native), and approaches
the energy floor because its reversible core is literally reversible. The
ladder is a roadmap, not a menu: one climbs it as fabrication allows, and the
software (\SYNARC{}) is unchanged at every rung because the logical
architecture is fixed (T-153).
